%==============================================================
\chapter{Methods}
\label{ch:methods}
%==============================================================
\section{Ultrafast Pulse Shaping}
\subsection{Ultraslow Centrifuge Design}

\begin{figure}[ht]
    \centering
    % ======================================================================
%  Ultraslow optical centrifuge pulse shaper (compressor arm + delay arm)
%  Included from chapters/03_methods.tex via \input.
% ======================================================================
\begin{tikzpicture}[line cap=butt,line join=round,>=Stealth,
  every node/.append style={font=\large\bfseries\boldmath}]

% =============================== nodes ================================
\coordinate (IN)   at (1.15,8.70);
\coordinate (HWP)  at (3.30,8.70);
\coordinate (PBSA) at (4.90,8.70);
\coordinate (IM)   at (7.465,8.70);
\coordinate (GRa)  at (4.36,12.02);
\coordinate (GRb)  at (10.75,12.05);
\coordinate (FM)   at (8.17,13.62);
\coordinate (RR)   at (12.65,13.62);
\coordinate (RRin) at (12.42,13.62);
\coordinate (OM)   at (9.85,6.15);
\coordinate (MG)   at (6.60,6.15);
\coordinate (PBSB) at (6.60,4.85);
\coordinate (QWP)  at (6.60,4.05);
\coordinate (MC)   at (4.90,5.70);
\coordinate (DA)   at (1.70,5.70);
\coordinate (DB)   at (1.70,4.85);

% ========================= translation stages =========================
\fill[stagegreen] (7.15,11.25) rectangle (11.55,14.45);
\fill[stageblue]  (0.95,4.20)  rectangle (3.15,6.35);

% =============================== beams ================================
% input pulse (purple)
\draw[cpurple,line width=3pt] (IN) -- (PBSA);
\fill[cpurple] (1.15,9.12) -- (1.15,8.28) -- (1.80,8.70) -- cycle;

% compressor arm (green)
\draw[cgreen,line width=3pt] (PBSA) -- (IM);
\draw[cgreen,line width=3pt] (GRa) -- (OM) -- (MG) -- (PBSB);

% delay arm (cyan)
\draw[ccyan,line width=3pt] (PBSA) -- (MC) -- (DA) -- (DB) -- (PBSB);

% output (purple)
\draw[cpurple,line width=3pt] (PBSB) -- (6.60,3.95);

% dispersed beams inside the compressor
% GR1 -> GR2 : angularly dispersed fan (red diffracted at the smaller angle, as in the inset)
\optbeam{GRa}{GRb}{0.05}{0.30}{cred}{cblue}
% GR2 -> FM  : collimated, spectrum spread transversally (order flipped by GR2)
\optbeam{GRb}{FM}{0.30}{0.30}{cblue}{cred}
% FM  -> RR  : order flipped again by FM
\optbeam{FM}{RRin}{0.30}{0.30}{cred}{cblue}

% ============================ components ==============================
% gratings: second argument = direction of the grating normal
\optgrating{GRa}{15}{1.25}    % Normale 15 deg  -> theta_0 = 15 deg zum Strahl GR1->GR2
\optgrating{GRb}{195}{1.25}   % Normale 195 deg -> GR2 exakt parallel zu GR1
% mirrors: second argument = direction of the mirror normal
\optmirror{FM}{-15.8}{1.00}   % 148.4 deg -> 0 deg
\optmirror{IM}{156.6}{1.00}   % 0 deg -> 133.1 deg (auf der Achse GR1--OM)
\optmirror{OM}{157}{1.00}     % -46.9 deg -> 180 deg
\optmirror{MG}{315}{0.90}     % 180 deg -> 270 deg
\optmirror{MC}{135}{0.90}     % 270 deg -> 180 deg
\optmirror{DA}{315}{0.90}     % 180 deg -> 270 deg
\optmirror{DB}{45}{0.90}      % 270 deg -> 0 deg

% retro-reflector (roof mirror, displaces the beam out of the drawing plane)
\begin{scope}[shift={(RR)}]
  \fill[black] (-0.22,0.46) -- (0.07,0) -- (-0.22,-0.46) --
               (0.22,-0.46) -- (0.22,0.46) -- cycle;
\end{scope}

% beam splitters and wave plates
\optpbs{PBSA}{0.72}
\optpbs{PBSB}{0.72}
\filldraw[fill=plateA,draw=black,line width=0.3pt] (3.21,8.28) rectangle (3.39,9.12);
\filldraw[fill=plateB,draw=black,line width=0.3pt] (6.18,3.96) rectangle (7.02,4.14);

% ========================== centrifuge field ==========================
% rotating linear polarisation (corkscrew), drawn as a twisted ribbon
% with a slight perspective tilt of the transverse axis
\foreach \i in {0,...,199}{%
  \pgfmathsetmacro{\sa}{\i/200}
  \pgfmathsetmacro{\sb}{(\i+1)/200}
  \pgfmathsetmacro{\pha}{760*pow(\sa,1.45)}
  \pgfmathsetmacro{\phb}{760*pow(\sb,1.45)}
  \pgfmathsetmacro{\Aa}{0.05+0.47*pow(\sa,1.15)}
  \pgfmathsetmacro{\Ab}{0.05+0.47*pow(\sb,1.15)}
  \pgfmathsetmacro{\ya}{3.95-3.15*\sa}
  \pgfmathsetmacro{\yb}{3.95-3.15*\sb}
  \pgfmathsetmacro{\xa}{6.60+\Aa*cos(\pha)}   \pgfmathsetmacro{\za}{\ya+0.30*\Aa*sin(\pha)}
  \pgfmathsetmacro{\xb}{6.60+\Ab*cos(\phb)}   \pgfmathsetmacro{\zb}{\yb+0.30*\Ab*sin(\phb)}
  \pgfmathsetmacro{\xc}{6.60-\Ab*cos(\phb)}   \pgfmathsetmacro{\zc}{\yb-0.30*\Ab*sin(\phb)}
  \pgfmathsetmacro{\xd}{6.60-\Aa*cos(\pha)}   \pgfmathsetmacro{\zd}{\ya-0.30*\Aa*sin(\pha)}
  \pgfmathsetmacro{\tone}{50+45*abs(sin(\pha))}
  \pgfmathsetmacro{\front}{sin(\pha)}
  \ifdim\front pt>0pt
    \filldraw[cpink!\tone!cviolet,line width=0.2pt]
      (\xa,\za) -- (\xb,\zb) -- (\xc,\zc) -- (\xd,\zd) -- cycle;
  \else
    \filldraw[cblue!\tone!cviolet,line width=0.2pt]
      (\xa,\za) -- (\xb,\zb) -- (\xc,\zc) -- (\xd,\zd) -- cycle;
  \fi}

% ========================= dimensions & labels ========================
\draw[<->,line width=0.9pt] (4.36,14.95) -- (10.75,14.95)
  node[midway,above=1pt]{$L$};
\draw[<->,line width=0.9pt] (1.90,7.15) -- (4.90,7.15)
  node[midway,above=1pt]{$\Delta t$};

\node at (3.95,13.00) {$\mathrm{GR}_1$};
\node at (10.95,10.85) {$\mathrm{GR}_2$};
\node at (12.65,12.70) {$\mathrm{RR}$};
\node[anchor=west] at (7.75,9.05) {$\mathrm{IM}$};
\node[anchor=west] at (10.35,6.55) {$\mathrm{OM}$};
\node[anchor=south west] at (5.20,9.15) {$\mathrm{PBS}$};
\node[anchor=west] at (7.05,4.85) {$\mathrm{PBS}$};
\node[plateA,anchor=south] at (3.30,9.22) {$\lambda/2$};
\node[plateB,anchor=east]  at (6.06,4.05) {$\lambda/4$};

% ============================== inset =================================
% definition of the grating orientation angle theta_0
\begin{scope}
  \fill[white] (8.60,0.55) rectangle (13.55,4.55);
  \draw[black,line width=1pt] (8.60,0.55) rectangle (13.55,4.55);
  \coordinate (GI) at (10.05,3.55);            % grating
  \coordinate (HI) at (13.25,3.55);            % direction of the outgoing beam
  \coordinate (NI) at ($(GI)+(-55:1.95)$);     % grating normal
  \coordinate (SI) at (13.05,4.28);            % source of the incident beam
  \draw[cgreen,line width=2.6pt] (SI) -- (GI);
  \optbeam{GI}{HI}{0.03}{0.34}{cblue}{cred}
  \draw[dashed,line width=0.7pt] (GI) -- (HI);
  \draw[dashed,line width=0.7pt] (GI) -- (NI);
  \draw[line width=0.6pt] ($(GI)+(0:0.80)$) arc (0:-55:0.80);
  \node at ($(GI)+(-28:1.10)$) {$\theta_0$};
  \optgrating{GI}{-55}{1.30}
\end{scope}

\end{tikzpicture}

    \caption{Schematic of the ultraslow centrifuge pulse shaper. The optical components used are: optical gratings (GR), polarizing beam splitters (PBS), mirrors (IM/ OM are the input and output mirrors of the compressor arm), retro-reflector (RR) and a half-wave plate (HWP) as well as a quarter-wave plate (QWP).}
    \label{fig:usCFG}
\end{figure}

As already stated in section \ref{sec:different_centrifuge_types} our new concept of the slow accelerating centrifuge is based on a Michelson interferometer.

To create the linear polarization rotating and accelerating in time, we use the uncompressed output of our chirped pulse amplifier (Spectra-Physics Spitfire Ace). Knowing the pulse dimensions  $\Delta \omega_{\mathrm{FWHM}} \approx 8.2 \, \mathrm{nm}$, central wavelength of $\lambda_0 = 802\, \mathrm{nm}$ and the duration $T_{\mathrm{FWHM}} \approx 320 \, \mathrm{ps}$, one can equate the temporal chirp with formula \ref{eq:chirp_from_widths} to $\left|\beta_0 \right| \approx 7.4 \times 10^{-2} \, \mathrm{rad}/\mathrm{ps}^2$. In our case positively chirped, meaning an increase of optical frequency with time. 

Looking at fig. \ref{fig:usCFG}, the input pulse first has its linear polarization rotated by a half-wave plate (HWP), so that the polarizing beam splitter (PBS) splits it into a horizontally and a vertically polarized arm of the interferometer. The HWP angle is not set to give an exact $50/50$ split at the PBS, but rather to balance the intensities of the two arms after the interferometer, compensating for the diffraction losses of the compressor arm. Unequal amplitudes at recombination would leave the centrifuge field elliptical rather than linear.

The transmitted beam (compressor arm, green in fig. \ref{fig:usCFG}) hits the first grating, which angularly disperses the spectral components. The second grating removes this angular dispersion, leaving a collimated but spatially chirped beam, and the retro-reflector sends it back through the pair at a different height, cancelling the spatial chirp. Since the red components are diffracted at larger angles, they traverse a longer optical path than the blue ones, so the grating pair adds negative group delay dispersion to the already positively chirped pulse. This shortens the pulse and thereby increases its chirp rate $\beta$. For a small difference in chirp, $|\Delta\beta| \ll |\beta_0|$, one can use the following approximation \footnote{Sign and factor conventions Ref. \cite{wangUltraslowOpticalCentrifuge2026} defines the temporal phase as $\varphi(t,\beta_0) = \varphi_0 + \omega_0 t + \beta_0 t^2$, i.e. their chirp coefficient is half the one used in this thesis, which accounts for the different prefactor (8 instead of 16).} \cite{wangUltraslowOpticalCentrifuge2026}:
\begin{equation}
     \Delta\beta(L) \approx \frac{8\beta_0^2\pi^2 c}{d^2\omega_0^3\cos^2\theta_0}\,L 
\end{equation}
where $1/d = 1500\,\mathrm{grooves/mm}$ is the groove density of the gratings, $\theta_0$ is the grating orientation angle (ref figure \ref{fig:usCFG}), and $L$ is the grating separation along the path of the light beam. The separation can be varied with a translation stage carrying the second grating, such that the optical path length at $\omega_0$ remains unchanged.

The translation stage in the delay arm (cyan in fig. \ref{fig:usCFG}) introduces the path length difference $\Delta t$ between the two pulses. After recombining the two beams at the second PBS, a quarter-wave plate (QWP) with its fast axis at $45^\circ$ to the polarization axes of the PBS converts the two orthogonal linear polarizations into circular polarizations of opposite handedness. Their superposition is a linearly polarized field whose polarization vector rotates at half the instantaneous frequency difference between the two arms.

Because changing $L$ in theory does not affect $\Delta t$, the two stages act as independent knobs: $\Delta t$ sets the central rotation frequency $f_0$ and $\Delta\beta(L)$ sets the angular acceleration, so that $f_{\mathrm{cfg}}(t)$ eq. \ref{eq:f_cfg} is fully controlled. Even though figure \ref{fig:centrifuge_types} c) makes it look like the centrifuge is only accelerating, it is absolutely possible to decelerate the rotation of the field. For example setting $\Delta t = 0$ would give a cfg that is decelerating first, then changes its rotation direction and accelerates the other way around.

\subsubsection{Restrictions to the current setup}
Even though it sounds like we can basically set any acceleration, in reality one is restricted among other things by the movement range of the translation stages. 
\todo{add typical frequency ranges, say that decelerating is also possible by tuning central frequency}


\subsection{Phase Control and Stabilization}


\section{Creating the Probing Environment}
The experiments are carried out in a four-chamber vacuum apparatus designed and built by Ian MacPhail-Bartley and shown in Figure. .... Most of the information in this chapter is taken from his PhD thesis  and for a more in depth description see [[ianThesis]].  There are four different chambers labeled source, doping, differential, and detection, separated by skimmers and apertures and pumped individually, so that the pressure drops from $\SI{e-4}{\torr}$ in the source chamber to below \SI{e-8}{\torr} in the detection chamber during operation.

\textbf{Source Chamber:} The source chamber is where helium nanodroplets are formed in a supersonic expansion from a cryogenically cooled nozzle (details in section ...). Only a small fraction of the expanding gas condenses into droplets, so a turbomolecular pump with a high pumping speed is attached to keep the chamber at $\sim\SI{2e-4}{\torr}$, corresponding to a mean free path of roughly \SI{50}{\centi\metre}. This ensures that droplets are not destroyed by collisions with background gas before they leave the chamber. A skimmer separates the source chamber from the doping chamber: it maintains the pressure difference between the two stages, removes the excess helium that has not condensed, and transmits only the central part of the expansion. Since condensation continues downstream of the nozzle for approximately 1000 nozzle diameters, the skimmer must be placed beyond this distance. If the gas density at the skimmer tip becomes too high, the skimmer itself begins to interfere with the beam. Incoming droplets are scattered by orifice clogging and helium atoms reflected from the inner walls.

\textbf{Doping Chamber:} A gate valve between the source and doping chambers allows background measurements to be taken without the droplet beam. With the valve open, the droplets pass through the doping cell, where they pick up dopant molecules according to Poisson statistics; these molecules serve as the probes of superfluidity in the following experiments. The cell is connected to an all-metal precision leak valve, which sets the dopant pressure inside the cell and thereby the average number of dopants per droplet (see Sec. ...). The chamber is pumped such that its pressure stays an order of magnitude below the cell pressure, so that pickup occurs essentially only within the cell.

\textbf{Differential Chamber:} The doped beam then passes into the differential pumping stage. Since the beam is no longer a supersonic expansion at this point, no shock structure forms in front of the entrance and a flat aperture is sufficient to provide the additional pressure drop. A second gate valve is installed after the doping chamber, which makes it possible to distinguish between two contributions to the dopant background in the detection chamber: effusive molecules travelling along the beam axis through the apertures, and molecules back streaming through the turbo molecular pumps (only the case because the pumps are all backed through the source chamber pump). The pulsed valve for supersonic jet experiments is also mounted in this chamber and can be lowered onto the beam axis, allowing convenient switching between droplet and jet measurements. The skimmer at the exit of this chamber is retained, as it is required to form the jet gas stream from the pulsed nozzle.

\textbf{VMI Chamber:} Last but not least the beam enters the chamber where the actual experiment is conducted. Here the laser pulses act on the probes and information is extracted by the time of flight velocity map imaging setup which will be explained in detail in sec ...

\subsection{Supersonic Gas Jet}
A pulsed supersonic jet is used to provide an environment for a reference measurement on isolated molecules. The molecules are seeded at a few percent in helium, which both improves the cooling and suppresses the formation of clusters (van der Waals complexes) of the probe molecule itself. The valve is a Parker Series 9, 200 µm orifice operated at $5-10\, \mathrm{Hz}$ with a backing pressure of $50-120\, \mathrm{psi}$. Under these conditions the temperature of the molecules is estimated to be around $10\, \mathrm{K}$. Because the molecules in the jet are collision-free at the interaction region and free of any surrounding medium, their rotational dynamics solely depend on the temperature distribution. Measurements under the centrifuge field can be compared directly against the corresponding measurement in helium droplets. Differences between the two are then attributable to the unequal initial rotational state distribution and the interaction with the superfluid.
\begin{aside}{Molecular beams and supersonic cooling}
A molecular beam is produced by expanding gas from a reservoir at stagnation conditions $(p_0,  T_0)$ through an orifice of diameter $D$ into vacuum. The character of the beam depends on the ratio of the mean free path $\lambda_0$ in the reservoir to $D$. For $\lambda_0 \gg D$ the molecules leave essentially without colliding and an effusive beam is obtained, with a velocity distribution corresponding to $T_0$. For $D \gg \lambda_0$ the molecules collide many times during the expansion, and since no heat is exchanged with the walls, the gas can only accelerate at the expense of its own enthalpy. The result is that random thermal motion is converted into a common forward velocity and since the temperature is set by the mean kinetic energy of the residual random motion in the co-moving frame, and this random component is progressively transferred into the directed flow, the beam becomes very cold while moving very fast. 
Unlike translational cooling, which results from binary collisions, condensation into clusters requires three-body collisions, since a third partner is needed to carry away the binding energy released when two monomers combine. This clustering is later crucial to create the helium nanodroplets. \cite{morse2SupersonicBeam1996}
\end{aside}
\subsection{Superfluid Helium Nanodroplets}
The clustering process introduced above becomes the dominant effect when helium itself is expanded. Two formation regimes are distinguished by whether the helium leaves the nozzle as a gas or as a liquid \cite{toenniesSuperfluidHeliumDroplets2004a}. If the stagnation state lies in the gaseous region, the expansion drives the gas past the vapor-liquid coexistence line and droplets nucleate and grow in the free jet, yielding $\langle N_{\mathrm{He}} \rangle \approx 10^2-10^5$ atoms. If instead $p_0$ and $T_0$ are such that the helium liquefies inside the nozzle, a liquid jet emerges and breaks up hydrodynamically, producing droplets of $10^6$ up to $10^{12}-10^{13}$ atoms.
\begin{figure}[ht]
    \centering
    \includegraphics[width=0.8\textwidth]{03_methods/droplets.png}
    \caption{}
    \label{fig:droplets}
\end{figure}
The droplets used here are produced in the gas-condensation regime, by a continuous expansion of 30 bar helium through a \SI{5}{\micro\meter} aperture at a nozzle temperature of \SIrange{16}{18}{\kelvin}, giving $\langle N_{\mathrm{He}} \rangle \approx 3000$ atoms \cite{ianThesis}.
Once formed, the droplets cool further in vacuum by evaporation of individual helium atoms, each removing approximately \SI{5}{\per\centi\meter} (\SI{7}{\kelvin}) of internal energy. The initial cooling rate is of order 10¹⁰ K/s and decays exponentially as the droplet cools, since the evaporation rate depends exponentially on temperature. After roughly 10⁻⁴ s — far shorter than the millisecond flight time through the apparatus — evaporation becomes negligible and the temperature stabilises at 0.37 K. The droplets therefore arrive at the interaction region well below the bulk λ-point of 2.17 K and in the superfluid state.

\subsubsection{Doping Nanodroplets with Molecules}

\section{Detection and Data Acquisition}
\label{sec:detection}
\todo{}

\section{Rotating Molecules: In-Field vs Field-free}
\label{sec:rotating_molecules}
\todo{spinnability? }
